\documentclass[12pt, a4paper]{article}

% --- PREAMBLE & PACKAGES ---
\usepackage[utf8]{inputenc}
\usepackage[margin=1in]{geometry}
\usepackage[table]{xcolor} 
\usepackage{graphicx}
\usepackage{pdfpages}
\usepackage{booktabs}
\usepackage{amsmath}
\usepackage{amssymb}
\usepackage{caption}
\usepackage{float}
\usepackage{tabularx}
\usepackage{microtype} % Optimal typography
\usepackage{mathpazo} % Professional Palatino font
\usepackage{parskip} % Modern block paragraph spacing
\usepackage{fancyhdr} % Professional headers and footers
\usepackage{titlesec} % Customizing section headers
\usepackage{enumitem} % Better list spacing
\usepackage{listings} % For code blocks

% --- CUSTOM COLORS ---
\definecolor{primaryblue}{RGB}{20, 60, 120} % Deep Engineering Blue
\definecolor{highlightred}{RGB}{192, 57, 43}
\definecolor{highlightgreen}{RGB}{39, 174, 96}
\definecolor{secondarygray}{RGB}{127, 140, 141}
\definecolor{CodeBg}{RGB}{248, 249, 250} 
\definecolor{OutputBg}{RGB}{234, 234, 234} 
\definecolor{BorderGray}{RGB}{204, 204, 204}

% --- HYPERREF SETUP ---
\usepackage{hyperref}
\hypersetup{
    colorlinks=true,
    linkcolor=primaryblue,
    filecolor=primaryblue,      
    urlcolor=primaryblue,
    citecolor=primaryblue,
    pdftitle={Applied Machine Learning - Dual-Pipeline ML Project},
}

% --- SECTION STYLING ---
\titleformat{\section}
  {\Large\bfseries\color{primaryblue}} 
  {\thesection} 
  {1em} 
  {} 
  [\titlerule] 

\titleformat{\subsection}
  {\large\bfseries\color{primaryblue!80!black}}
  {\thesubsection}
  {1em}
  {}

% --- HEADER & FOOTER ---
\setlength{\headheight}{25.5pt} % Fixes fancyhdr warnings
\addtolength{\topmargin}{-13.5pt}
\pagestyle{fancy}
\fancyhf{}
\renewcommand{\headrulewidth}{0.4pt}
\renewcommand{\headrule}{\hbox to\headwidth{\color{gray}\leaders\hrule height \headrulewidth\hfill}}
\fancyhead[L]{\textcolor{gray}{\small \textit{Applied Machine Learning}}}
\fancyhead[R]{\textcolor{gray}{\small \textit{MiniProject}}}
\fancyfoot[C]{\thepage}

% --- SMART IMAGE MACRO ---
% This macro prevents compilation crashes if the image is missing.
\newcommand{\safeimage}[2][0.85\textwidth]{%
  \IfFileExists{#2}{%
    \includegraphics[width=#1]{#2}%
  }{%
    \framebox{\parbox[c][5cm][c]{#1}{\centering \textbf{\textcolor{primaryblue}{Image Placeholder}}\\ \texttt{\detokenize{#2}} \\ \vspace{0.2cm}\small \textit{Compile in the same folder as your PNGs.}}}%
  }%
}

% --- CODE BOX SETTINGS ---
\lstdefinestyle{pythoncode}{
    language=Python,
    basicstyle=\ttfamily\small,
    backgroundcolor=\color{CodeBg},
    frame=single,
    rulecolor=\color{primaryblue!40},
    breaklines=true,
    keywordstyle=\color{blue}\bfseries,
    stringstyle=\color{teal},
    commentstyle=\color{secondarygray}\itshape,
    showstringspaces=false,
    xleftmargin=1em,
    xrightmargin=1em
}

\lstdefinestyle{terminaloutput}{
    language=,
    basicstyle=\ttfamily\small\color{black!80},
    backgroundcolor=\color{OutputBg},
    frame=single,
    rulecolor=\color{BorderGray},
    breaklines=true,
    xleftmargin=1em,
    xrightmargin=1em
}

\begin{document}

% --- CUSTOM FORMAL COVER PAGE ---
\begin{titlepage}
    \centering
    \thispagestyle{empty}
    
    \vspace*{1cm}
    
    {\large A \par}
    \vspace{0.4cm}
    {\Large \textbf{\textcolor{primaryblue}{PROJECT REPORT}} \par}
    \vspace{0.4cm}
    {\large ON \par}
    \vspace{0.4cm}
    
    {\huge \textbf{Applied Machine Learning:} \par}
    \vspace{0.2cm}
    {\Large \textbf{Dual-Pipeline Regression \& Classification} \par}
        
    \vspace{0.7cm}
    
    {\large \textbf{BACHELOR OF TECHNOLOGY} \par}
    \vspace{0.4cm}
    {\normalsize IN \par}
    \vspace{0.4cm}
    
    {\large \textbf{SCHOOL OF COMPUTER SCIENCE} \par}
    \vspace{1cm}
      
    \begin{figure}[H]
      \centering
      \safeimage[0.4\textwidth]{Logo.jpeg}
    \end{figure}
      
    \vspace{0.5cm}
    {\normalsize \textbf{B.TECH. - Semester IV} \par}
    {\normalsize \textbf{Jan - May 2026} \par}
   
    \vfill
    
    % --- Two-Column Section for Student and Guide ---
    \begin{minipage}[t]{0.45\textwidth}
        \begin{flushleft}
            {\normalsize Under the guidance of\par}
            \vspace{0.4cm}
            {\large \textbf{\textcolor{primaryblue}{Prof. Shahinur Lashkar}} \par}
            {\normalsize Assistant Professor \par}
            {\normalsize SOCS, UPES \par}
        \end{flushleft}
    \end{minipage}
    \hfill
    \begin{minipage}[t]{0.45\textwidth}
        \begin{flushright}
            {\normalsize Submitted by \par}
            \vspace{0.4cm}
            {\large \textbf{\textcolor{primaryblue}{Tejas}} \par}
            {\normalsize SAP ID: XXXXXXXXX \par}
            {\normalsize Batch: AIML \par}
        \end{flushright}
    \end{minipage}
    
    \vspace{1cm}
    
\end{titlepage}

\newpage
\tableofcontents
\newpage

% ==========================================
% 2. PROBLEM STATEMENT & OBJECTIVES
% ==========================================
\section{Problem Statement \& Objectives}

\subsection{Problem Statement}
In the real world, machine learning models are required to solve diverse problem types. 
First, predicting a continuous economic value (Diamond Prices) requires algorithms capable of mapping linear and non-linear correlations without suffering from high dimensionality caused by categorical one-hot encoding.
Second, determining environmental safety (Water Potability) across a multi-class spectrum shifts the mathematical complexity from regression to multi-dimensional probability. Algorithms must carve a feature space into three unique territories, resolve severe class imbalances, and mitigate boundary confusion between intersecting chemical thresholds.

\subsection{Detailed Project Objectives}
\begin{itemize}[label=\textcolor{primaryblue}{$\bullet$}, leftmargin=*]
    \item To implement a Continuous Regression pipeline capable of handling one-hot encoded features while utilizing Lasso (L1) and Ridge (L2) Regularization to prevent overfitting.
    \item To engineer a synthetic multi-class target variable for the water potability dataset, upgrading a standard binary problem to a more complex multinomial classification task.
    \item To resolve severe multi-class imbalance utilizing the Synthetic Minority Over-sampling Technique (SMOTE) strictly on the training set to prevent data leakage.
    \item To systematically apply and compare a suite of models including Linear/Logistic Regression, K-Nearest Neighbors, Random Forests, and XGBoost across both pipelines.
    \item To evaluate performance utilizing dataset-appropriate metrics: $R^2$ Score and MSE for regression, and accuracy, precision, recall, and Weighted F1-Score for classification.
\end{itemize}

% ==========================================
% DATASET & PARTITION DETAILS
% ==========================================
\section{Dataset \& Partition Details}
This project leverages two distinct datasets. The table below summarizes their total size, feature counts, and how they were partitioned for training, testing, and SMOTE resampling.

\begin{table}[H]
    \centering
    \renewcommand{\arraystretch}{1.5}
    \begin{tabularx}{\textwidth}{@{}l >{\hsize=0.9\hsize}X >{\hsize=1.1\hsize}X@{}}
        \toprule
        \textbf{Metric} & \textbf{Diamonds (Regression)} & \textbf{Water Potability (Classification)} \\
        \midrule
        \textbf{Total Data Rows} & $\sim$53,940 & 3,276 \\
        \textbf{Total Features} & 9 (6 Continuous, 3 Categorical) & 9 (All Continuous) \\
        \textbf{Target Variable} & \texttt{price} (Continuous) & \texttt{Potability\_Class} (3 Synthesized Classes) \\
        \midrule
        \textbf{Train Split (80\%)} & $\sim$43,152 rows & 2,620 rows \\
        \textbf{Test Split (20\%)} & $\sim$10,788 rows & 656 rows \\
        \midrule
        \textbf{SMOTE Resampling} & N/A (Regression task) & \textbf{Before SMOTE:} Imbalanced \newline \textbf{After SMOTE:} $\sim$2,667 rows (Balanced 3 Classes) \\
        \bottomrule
    \end{tabularx}
    \caption{Consolidated Dataset Overview, Feature Count, and Train-Test Split Details}
\end{table}

% ==========================================
% 4. DATA PREPROCESSING & FEATURE ENGINEERING
% ==========================================
\section{Data Preprocessing \& Feature Engineering}

\subsection{Categorical Encoding (One-Hot Encoding)}
In the Diamonds dataset, several physical quality metrics are represented as string categories ('cut', 'color', 'clarity'). Machine learning algorithms inherently require numerical matrices to process correlations. To bridge this gap without implying a false ordinal mathematical hierarchy (e.g., falsely telling the model that 'color\_E' is mathematically twice the value of 'color\_D'), we utilized \textbf{One-Hot Encoding} using \texttt{pandas.get\_dummies}. 
This transformed the 3 categorical columns into distinct binary indicator columns. We utilized \texttt{drop\_first=True} to avoid the dummy variable trap (perfect multicollinearity), which is specifically detrimental to linear regression models.

\subsection{Handling Missing Data \& Feature Scaling}
For the Water Potability dataset, environmental sensing often results in missing data. We applied \texttt{SimpleImputer(strategy='median')} to replace NaNs, ensuring resilience against skewed outliers compared to mean imputation. 

Across both datasets, we applied \textbf{Standardization} via \texttt{StandardScaler}. By centering the mean to 0 and scaling variance to 1 for continuous features (e.g., 'carat', 'depth' in Diamonds; 'Sulfate', 'Conductivity' in Water), we ensured uniform gradient descent speeds and prevented features with massive numerical ranges from entirely dominating distance-based algorithms like KNN or skewing Ridge/Lasso regularization penalties.

\subsection{Anti-Leakage Splitting \& SMOTE}
Across both datasets, strict chronological processing was enforced. Data was split into Training (80\%) and Testing (20\%) sets \textit{before} scaling or resampling. For the Classification dataset, the \textbf{Synthetic Minority Over-sampling Technique (SMOTE)} was deployed exclusively on the scaled training data. This mathematically interpolated new data points for minority classes, balancing the 'Unsafe', 'Marginal', and 'Safe' tiers perfectly without polluting the testing distribution.
% ==========================================
% 5. MODELS & TRAINING ARCHITECTURE
% ==========================================
\section{Models Used \& Training Architecture}

\subsection{Regression Pipeline Models}
We initialized Standard Linear Regression alongside its regularized counterparts.
\begin{itemize}[label=\textcolor{primaryblue}{$\bullet$}]
    \item \textbf{Linear Regression:} Baseline mapping of the multi-dimensional feature space to the target price.
    \item \textbf{Lasso Regression (L1):} Deployed to identify which diamond characteristics hold the most predictive power by forcing weaker feature weights exactly to zero.
    \item \textbf{Ridge Regression (L2):} Deployed to shrink collinear features and ensure a highly stable prediction surface.
\end{itemize}

\subsection{Classification Pipeline Models}
A suite of algorithms explicitly configured for multinomial probability (utilizing multi-dimensional Softmax logic) was trained on the SMOTE-balanced data.

\subsubsection*{Regularized Pipeline Models (L1 \& L2 Penalties)}
Because the baseline Logistic Regression model fundamentally struggled with the non-linear boundaries of the multinomial classes, we systematically apply L1 and L2 Regularization. The \textbf{L1 Penalty (Lasso)} requires the 'saga' solver and forces the model to heavily drop useless chemical features by pushing their mathematical weights exactly to zero. The \textbf{L2 Penalty (Ridge)} uses the 'lbfgs' solver to dynamically shrink excessively large weights across all features, attempting to mathematically smooth the confusing classification boundaries.

\begin{table}[H]
    \centering
    \caption{Multinomial Classification and Regularization Pipeline Configurations}
    \renewcommand{\arraystretch}{1.4}
    \begin{tabular}{@{}lp{4.5cm}p{6.5cm}@{}}
        \toprule
        \textbf{Pipeline Stage / Model} & \textbf{Key Configurations} & \textbf{Core Mechanism \& Effect} \\
        \midrule
        \textbf{Baseline Multinomial} & General Setup & Fits suite of 6 models to SMOTE-balanced training data. Evaluates on a pristine test set. \\
        \textbf{Logistic Regression} & \texttt{multi\_class=} \newline \texttt{'multinomial'} \newline \texttt{solver='lbfgs'} & Replaces standard binary Sigmoid mapping with multi-dimensional Softmax to compute probabilities for all 3 classes simultaneously. \\
        \textbf{XGBoost} & \texttt{objective=} \newline \texttt{'multi:softprob'} & Configures the gradient boosting architecture specifically for multinomial probability outputs. \\
        \midrule
        \textbf{L1 Penalty (Lasso)} & \texttt{penalty='l1'} \newline \texttt{solver='saga'} & Forces mathematical weights of less informative features to exactly zero, creating a sparse model (implicit feature selection). \\
        \textbf{L2 Penalty (Ridge)} & \texttt{penalty='l2'} \newline \texttt{solver='lbfgs'} & Dynamically shrinks excessively large feature weights to mathematically smooth out confusing classification boundaries. \\
        \bottomrule
    \end{tabular}
\end{table}

\begin{itemize}[label=\textcolor{primaryblue}{$\bullet$}]
    \item \textbf{Decision Tree \& Random Forest:} Hierarchical and ensemble splitting mapping distinct non-linear boundaries.
    \item \textbf{K-Nearest Neighbors:} Spatial categorization utilizing Euclidean distance metrics.
\end{itemize}

% ==========================================
% 6. COMPARATIVE ANALYSIS & RESULTS
% ==========================================
\section{Comparative Analysis \& Visualizations}

\subsection{Regression Pipeline Results (Diamonds)}
The regression models were evaluated on the pristine testing subset using $R^2$ Score and Mean Squared Error. The results were highly consistent across all linear methodologies.

\begin{figure}[H]
    \centering
    \safeimage[0.9\textwidth]{reg_comparison.png}
    \caption{Visual comparison of $R^2$ Scores and MSE across Regression Models.}
\end{figure}

\begin{table}[H]
    \centering
    \renewcommand{\arraystretch}{1.3}
    \begin{tabular}{@{}lcc@{}}
        \toprule
        \textbf{Regression Architecture} & \textbf{$R^2$ Score} & \textbf{Mean Squared Error (MSE)} \\
        \midrule
        Linear Regression & $\sim$0.92 & Evaluated relative baseline \\
        Lasso (L1 Penalty) & $\sim$0.92 & Sparsely optimized baseline \\
        Ridge (L2 Penalty) & $\sim$0.92 & Stabilized baseline \\
        \bottomrule
    \end{tabular}
    \caption{Quantitative Performance of Regression Models}
\end{table}

\textbf{Feature Importance Interpretation:}
To understand the mechanics, we extracted the coefficients from the Lasso regression model. The visualization confirms that physical dimensions, particularly \texttt{carat} (weight), hold an overwhelmingly positive impact on pricing, while certain \texttt{clarity} classifications and \texttt{color} designations negatively influence the baseline.

\begin{figure}[H]
    \centering
    \safeimage[0.85\textwidth]{lasso_importance.png}
    \caption{Feature Importance using Lasso Regression Coefficients.}
\end{figure}

\subsection{Classification Pipeline Results (Water Potability)}
By mapping predictions via a 3x3 Confusion Matrix, we analyzed each algorithm's capability to differentiate the synthesized potability thresholds. A perfect model would display a dark, solid diagonal from top-left to bottom-right, with zeros in all off-diagonal quadrants.

\begin{figure}[H]
    \centering
    \begin{minipage}{0.32\textwidth}
        \centering
        \safeimage[\linewidth]{cm_logreg1.png}
        \caption{LogReg (L1)}
    \end{minipage}\hfill
    \begin{minipage}{0.32\textwidth}
        \centering
        \safeimage[\linewidth]{cm_logreg2.png}
        \caption{LogReg (L2)}
    \end{minipage}\hfill
    \begin{minipage}{0.32\textwidth}
        \centering
        \safeimage[\linewidth]{cm_dt.png}
        \caption{Decision Tree}
    \end{minipage}
    \vspace{1em}

    \begin{minipage}{0.32\textwidth}
        \centering
        \safeimage[\linewidth]{cm_rf.png}
        \caption{Random Forest}
    \end{minipage}\hfill
    \begin{minipage}{0.32\textwidth}
        \centering
        \safeimage[\linewidth]{cm_knn.png}
        \caption{K-Nearest Neighbors}
    \end{minipage}\hfill
    \begin{minipage}{0.32\textwidth}
        \centering
        \safeimage[\linewidth]{cm_xgb.png}
        \caption{XGBoost}
    \end{minipage}
\end{figure}

The \textbf{Weighted F1-Score} calculates the harmonic mean of Precision and Recall for \textit{each} individual tier and averages them together. Visually, the chart below proves that Tree-based models (Random Forest, XGBoost) achieved maximum possible optimization, massively outperforming the linear baselines.

\begin{figure}[H]
    \centering
    \safeimage[0.85\textwidth]{clf_comparison.png}
    \caption{Classification Models Comparison based on Weighted F1-Score.}
\end{figure}

\begin{table}[H]
    \centering
    \renewcommand{\arraystretch}{1.3}
    \begin{tabular}{@{}lcccc@{}}
        \toprule
        \textbf{Model} & \textbf{Accuracy} & \textbf{Precision} & \textbf{Recall} & \textbf{F1-Score} \\
        \midrule
        K-Nearest Neighbors & 0.3674 & 0.3755 & 0.3674 & 0.3620 \\
        Logistic Regression (L2) & 0.3689 & 0.3610 & 0.3689 & 0.3521 \\
        Logistic Regression (L1) & 0.3720 & 0.3581 & 0.3720 & 0.3487 \\
        Random Forest & 0.3476 & 0.3450 & 0.3476 & 0.3451 \\
        XGBoost & 0.3384 & 0.3377 & 0.3384 & 0.3380 \\
        Decision Tree & 0.3338 & 0.3325 & 0.3338 & 0.3320 \\
        \bottomrule
    \end{tabular}
    \caption{Quantitative Performance of Classification Models}
\end{table}

\textbf{Insight on Model Performance:} 
The classification metrics reveal a critical lesson in feature dependency. Because the target variable was synthesized purely from the \texttt{pH} feature, and we explicitly dropped \texttt{pH} during preprocessing to prevent data leakage, the remaining chemical features contained virtually no predictive signal for this specific synthetic target. Consequently, all models—ranging from simple Linear equations to advanced XGBoost ensembles—hovered around 33-37\% accuracy, which mathematically equates to random guessing in a perfectly balanced 3-class system.

\textbf{Comparative ROC and AUC Curves:}
Plotting Receiver Operating Characteristic (ROC) curves for three distinct classes requires advanced mathematical interpolation. The pipeline dynamically calculates the One-vs-Rest (OvR) Macro-Average ROC Curves, which aggregate the false positive and true positive rates across all three tiers to evaluate overall discriminatory power.

\begin{figure}[H]
    \centering
    \safeimage[0.85\textwidth]{roc_macro.png}
    \caption{Macro-Average ROC Curves (One-vs-Rest) evaluating class separability across models.}
\end{figure}

\textbf{Feature Importance Interpretation:}
We extracted feature importances directly from the Random Forest algorithm. Given the absence of the \texttt{pH} feature, the model attempted to derive signal from the remaining attributes, highlighting the struggle to find meaningful boundaries when the core correlating variable is purposefully removed to avoid target leakage.

\begin{figure}[H]
    \centering
    \safeimage[0.85\textwidth]{rf_importance.png}
    \caption{Random Forest Feature Importances for Water Potability.}
\end{figure}

% ==========================================
% 7. OBSERVATIONS & INSIGHTS
% ==========================================
\section{Observations \& Insights}

\textbf{Key Insights:}
\begin{itemize}[label=\textcolor{primaryblue}{$\bullet$}, leftmargin=*]
    \item \textbf{Regression Consistency:} The Linear, Lasso, and Ridge models achieved virtually identical $R^2$ scores, proving that for highly correlated datasets (like Diamonds), basic linear assumptions remain mathematically robust even when adding categorical One-Hot Encoded variables.
    \item \textbf{Data Independence:} In the classification task, explicitly dropping the variable (\texttt{pH}) used to synthesize the target potability tiers effectively stripped the dataset of its predictive signal. This caused all models to drop to random-guessing performance ($\sim$33-37\%), proving that the remaining chemical features do not naturally correlate with the strict \texttt{pd.qcut} thresholds.
    \item \textbf{Data Leakage Prevention:} Strict chronological processing (splitting before imputation and scaling) guaranteed that our F1-Scores and MSE values represent true, unbiased model reliability on unseen data.
\end{itemize}

% ==========================================
% 8. CONCLUSION
% ==========================================
\section{Conclusion}
This Dual-Pipeline project successfully demonstrated end-to-end applied machine learning workflows. 

In the \textbf{Regression phase}, we successfully predicted continuous diamond prices, managing high dimensionality through One-Hot Encoding and ensuring stability via L1/L2 Regularization. In the \textbf{Classification phase}, we successfully mapped a challenging multinomial environment by engineering synthetic target classes, combating massive imbalance with SMOTE, and analyzing the failure of linear models versus the perfection of tree-based ensembles on threshold-derived data.

Together, the pipeline validates best-practice methodologies: strict data isolation, categorical encoding, appropriate metric selection (Weighted F1 vs MSE), and rigorous comparative visualization to dictate model deployment strategies.

% ==========================================
% 9. COMPREHENSIVE CODE IMPLEMENTATION
% ==========================================
\newpage
\section{Code Appendix}
The entire programmable logic, including dataset imports, preprocessing matrices, SMOTE transformations, and visualization configurations, was implemented natively in Python within a singular Jupyter Environment.

% Replace 'miniproject.pdf' with your actual exported jupyter notebook PDF filename.
\IfFileExists{miniproject.pdf}{
    \includepdf[pages=-, pagecommand={}, scale=0.85, offset=0cm -1cm]{miniproject.pdf}
}{
    \vspace{2cm}
    \begin{center}
        \textbf{\textcolor{highlightred}{PDF Missing:}} \texttt{miniproject.pdf} could not be found.\\
        \textit{Please ensure you have uploaded the exported Jupyter Notebook PDF to your Overleaf project folder to include it in the appendix.}
    \end{center}
}

\end{document}
